% Content-depth probe across three model families and five data splits.
\begin{table}[t]
\centering
\small
\setlength{\tabcolsep}{4pt}
\resizebox{\columnwidth}{!}{%
\begin{tabular}{@{}lcccc@{}}
\toprule
\textbf{Model} & \textbf{$L$} & \textbf{linear knee} & \textbf{native knee} & \textbf{gap} \\
 & & ($/L$, 95\% CI) & ($/L$) & ($\Delta/L$) \\
\midrule
Qwen3-8B & 36 & 0.393 [0.319,0.466] & 0.824 & 0.43 \\
Llama-3-8B$^\dagger$ & 32 & 0.275 & 1.000 & 0.725 \\
OLMo-2-7B & 32 & 0.285 [0.214,0.357] & 0.875 & 0.59 \\
\bottomrule
\end{tabular}%
}
\caption{\textbf{Linear accessibility precedes native readout under these
probes.} The linear knee is the first layer reaching 98\% of peak held-out
probe accuracy; the native knee applies each model's own final norm and LM head
at every layer. Qwen and OLMo aggregate SST2, WiC, and RTE over five stratified
data splits; their intervals reflect task--split variation, not model-training
seeds. $^\dagger$For Llama, WiC and RTE native verbalizers remain near chance,
so we report the SST2-only matched-support knees and do not form a misleading
three-task native aggregate. These correlational probes do not localize
understanding or establish the same knees on long-memory tasks.}
\label{tab:depth}
\end{table}
